
We began this work by asking: Why do Kaggle winners and PyTorch tutorials systematically avoid horizontal flip augmentation on MNIST digits? Our experiments provide a clear answer for MNIST: horizontal flip introduces label noise for asymmetric digits, degrading test accuracy by 0.72-1.00 percentage points at moderate flip rate (p=0.5), with dose-dependent degradation ranging from 0.37 pp (p=0.3) to 4.10 pp (p=0.9), while rotation $\pm 15$° causes no such harm. Practitioner intuition was correct for MNIST---but until now, lacked quantitative validation and mechanistic understanding.

The evidence establishes a clear causal chain on MNIST. Horizontal flip creates visually non-canonical images of asymmetric digits {2, 3, 5, 6, 7, 9}, yet retains their original labels, thereby injecting label noise proportional to flip probability. Models trained on this label-noisy data exhibit degraded accuracy on canonical test digits, with degradation magnitude increasing monotonically with flip rate. The perfect or near-perfect dose-response relationships ($\rho$ = -1.0 or -0.969) reveal a deterministic mechanism. Meanwhile, symmetric digits {0, 1, 8} remain largely stable at moderate flip rates (<0.2\% change at p=0.5), and rotation augmentation produces no differential degradation across four independent validations.

These findings demonstrate \textbf{semantic validity testing on MNIST} and propose a generalizable framework: an augmentation is semantically valid if and only if the augmented image remains in-distribution for its assigned label. Horizontal flip violates this criterion for asymmetric MNIST digits; rotation $\pm 15$° respects it. This principle is hypothesized to extend beyond MNIST to domains with semantic asymmetry---medical imaging, traffic signs, character recognition---though effect sizes and mechanisms require empirical testing on each dataset.

Our work opens several research directions. First, generalization: Does the semantic validity principle hold for Fashion-MNIST, CIFAR-10, and chest X-ray anatomical constraints? Second, architectural robustness: Do modern architectures or pre-trained models mitigate semantic invalidity through increased capacity? Third, augmentation scope: Does vertical flip harm digits 6 and 9, and can cutout or mixup introduce semantic violations? Finally, prescriptive solutions: Can semantic validity constraints be integrated into automated augmentation search?

\textbf{Semantic validity is not folklore---it is a testable design principle demonstrated on MNIST.} When augmentations violate domain-specific semantic constraints, they introduce label noise with predictable, quantifiable consequences. The perfect or near-perfect dose-response relationships we observed on MNIST indicate a deterministic mechanism rather than statistical artifact. Future augmentation policies should be grounded in principled validation: test dose-response relationships, employ positive controls, and measure class-specific effects rather than relying on aggregate metrics that mask selective harm.

The next time you design an augmentation policy for domains with semantic asymmetry, ask not ''What transformations are standard?'' but ''Do these transformations preserve class identity in my domain?'' For MNIST, the answer is clear: avoid horizontal flip, use rotation.

